% Generated draft: requires mathematical and source review.
\documentclass[11pt]{article}
\usepackage[margin=25mm]{geometry}
\usepackage{fontspec}
\setmainfont{Latin Modern Roman}
\usepackage{amsmath,amssymb,mathtools,unicode-math}
\setmathfont{Latin Modern Math}
\usepackage{enumitem,needspace,xurl,hyperref}
\setlength{\parindent}{0pt}
\setlength{\parskip}{5pt}
\setlength{\emergencystretch}{3em}
\newcommand{\N}{\mathbb{N}}
\newcommand{\Z}{\mathbb{Z}}
\newcommand{\Q}{\mathbb{Q}}
\newcommand{\R}{\mathbb{R}}
\newcommand{\C}{\mathbb{C}}
\newcommand{\F}{\mathbb{F}}
\newcommand{\A}{\mathbb{A}}
\newcommand{\PP}{\mathbb{P}}
\newcommand{\E}{\mathbb{E}}
\newcommand{\eps}{\varepsilon}
\newcommand{\dd}{\,\mathrm{d}}
\begin{document}
\section*{UnsolvedMath ID: 30005888}
\textbf{Source problem code:} catalogue entry
\subsection{Definitions and mathematical statement}
Let \(q\) be a prime power, let \(\F_q\) be the field with \(q\) elements, and write
\[
G=\operatorname{SL}_n(\F_q)=\{g\in\operatorname{Mat}_{n\times n}(\F_q):\det g=1\}.
\]
The supplied record asks for improved diameter bounds for these groups and gives the proposed scale
\[
O(1)(\log |G|)(\log n).
\]
The record does not specify the generating sets, the convention for diameter, the range of parameters, or the permitted dependence of the implicit constant. These omissions prevent an unconditional, uniquely quantified formulation.

For clarity, one possible convention is the following. For a generating subset \(S\subseteq G\), define the undirected word diameter by
\[
\operatorname{diam}(G,S)=\min\{\ell\in\Z_{\geq0}:G=(S\cup S^{-1}\cup\{1\})^\ell\}.
\]
Under this convention, if the intended question concerns all generating subsets with an absolute implicit constant, its candidate formulation is: does there exist a constant \(C>0\), independent of \(n,q,S\), such that for every integer \(n\geq2\), every prime power \(q\), and every generating subset \(S\subseteq\operatorname{SL}_n(\F_q)\),
\[
\operatorname{diam}(\operatorname{SL}_n(\F_q),S)
\leq C\log|\operatorname{SL}_n(\F_q)|\log n?
\]
This is an explicitly conditional interpretation of the supplied expression, not a claim that its quantifiers or diameter convention have been established from the source. Possible exceptional parameters or asymptotic restrictions are not supplied.

For a subset \(A\) of a group and an integer \(k\geq1\), the usual product-set notation is
\[
A^k=\{a_1\cdots a_k:a_1,\ldots,a_k\in A\}.
\]
The accompanying source question asks what can be said about the size and structure of the expression printed as \((A^k)(\Sigma)\). The cleaned statement describes subsets with ``permutation-rich restrictions'', but supplies no definition of that condition. The ambient group for \(A\), the set \(\Sigma\), the meaning of the parenthesized operation, the hypotheses on \(A\), the permissible exponent \(k\), and the desired size or structural conclusion are all unspecified. Consequently, this part cannot be turned into a precise quantified assertion from the supplied record alone; in particular, \((A^k)(\Sigma)\) is retained without identifying it with an image, restriction, orbit, or stabilizer construction.

The original-statement field also contains adjacent questions whose inclusion in this permanent record is uncertain. It reports that S. Eberhard posed four questions, three labeled (Q1), (Q2), and (Q3) in his extended abstract, and gives the fourth as Question 5: what is the length of the shortest law on \(S_n\)? Here a law conventionally means a nontrivial reduced word in a free group that evaluates to the identity under every substitution of elements of the group; its length is its reduced word length. The excerpt reports a lower bound printed as \(2n-1\) and an upper bound
\[
\exp\!\bigl(O((\log n)^4\log\log n)\bigr),
\]
with the upper bound attributed to Kozma--Thom, reference [15]. The typography of the reported lower bound may have lost mathematical formatting; no correction is inferred here, and neither bound is independently verified.

The same field attributes two further questions to M. Kassabov: Question 8 asks whether a random generating set in \(\operatorname{Alt}(n)\) gives expanders; Question 9 asks for which function \(f(n)\) there exists an expanding generating set in an expression rendered as \(\operatorname{Alt}(n)f(n)\). The sampling rule, number of generators, and uniform expansion requirements are unspecified. The latter expression may have lost exponent formatting, but a direct-power interpretation is not asserted. These neighboring questions are preserved as source context rather than merged into the diameter and product-structure target.
\subsection{Short English statement}
The record asks whether special linear groups admit diameter bounds of order \(\log|G|\log n\), and asks about the size and structure of an undefined product-set expression \((A^k)(\Sigma)\). The diameter formulation above is conditional on an explicit convention; the second question lacks essential definitions. Adjacent questions and their attributions are retained with their uncertainties. This is a statement-only draft requiring human mathematical review.
\subsection{Sources}
\begin{itemize}
\item Local UnsolvedMath registry, \texttt{data/unsolvedmath/problems.json}, supplied complete source record with permanent ID \(30005888\), source problem code \texttt{catalogue entry}, title ``Diameter and Product Structure in High-Rank Linear Groups'', and proposed year 2024. The supplied collection is ``Oberwolfach Reports Open Problems''. No individual proposer is specified for the principal target.
\item Supplied report citation: \emph{Mini-Workshop: Growth and Expansion in Groups} (2024), Oberwolfach Reports, pp.\ 1005--1040, DOI 10.4171/owr/2024/17. Supplied source link: \url{https://doi.org/10.4171/owr/2024/17}. The background describes Oberwolfach Reports as documenting mathematical workshops and their open-problem discussions; no precise page or question number within the report is supplied for the principal target.
\item Supplied literature source, labeled ``classical-group diameter bounds'': \url{https://arxiv.org/abs/2110.02942}. The record supplies no author, title, or theorem identification for this link.
\item Supplied literature source, labeled ``OWR report'': \url{https://ems.press/content/serial-article-files/49475}.
\item Literature assessment checked at \texttt{2026-08-22}: the supplied curation labels the proposed near-linear-log diameter bound and associated product-structure question open, and states that available classical-group bounds retain substantial rank dependence. This is supplied dated curation derived from the OWR material, not an independently verified assessment. The source explicitly requests independent verification of time-sensitive status and cited claims.
\item The original-statement excerpt attributes neighboring questions to S. Eberhard and M. Kassabov, and attributes the displayed upper bound for shortest laws on \(S_n\) to Kozma--Thom, reference [15]. No bibliographic entry for [15] is supplied.
\end{itemize}
\end{document}
